\documentclass[11pt]{article}

\usepackage[margin=1in]{geometry}
\usepackage{xcolor}
\usepackage{amsmath}
\usepackage{graphicx}
\usepackage{listings}
\usepackage{xurl}
\usepackage{hyperref}
\usepackage{booktabs}
\usepackage{enumitem}
\usepackage{parskip}
\usepackage{titlesec}
\usepackage{tcolorbox}
\tcbuselibrary{skins,breakable}

\hypersetup{
  colorlinks = true,
  linkcolor  = blue!60!black,
  urlcolor   = blue!60!black,
  citecolor  = blue!60!black,
}

% --- Code listing style -------------------------------------------------------
\definecolor{shellbg}{HTML}{F5F5F5}
\definecolor{shellrule}{HTML}{CCCCCC}
\definecolor{shellcomment}{HTML}{6A737D}
\definecolor{shellkeyword}{HTML}{005CC5}

\lstdefinestyle{shell}{
  basicstyle      = \ttfamily\footnotesize,
  backgroundcolor = \color{shellbg},
  frame           = single,
  rulecolor       = \color{shellrule},
  framesep        = 6pt,
  xleftmargin     = 6pt,
  xrightmargin    = 6pt,
  breaklines      = true,
  showstringspaces= false,
  columns         = fullflexible,
  keepspaces      = true,
  commentstyle    = \color{shellcomment}\itshape,
  keywordstyle    = \color{shellkeyword}\bfseries,
  morecomment     = [l]{\#},
  morekeywords    = {magic, export, source, cd, git, klayout, irsim, netgen,
                     sudo, apt-get, dnf},
  literate        = {/}{{/}}1 {-}{{-}}1 {=}{{=}}1,
}
\lstset{style=shell}

% --- Tip / Warning boxes ------------------------------------------------------
\newtcolorbox{tipbox}{
  colback  = blue!5!white,
  colframe = blue!60!black,
  fonttitle= \bfseries,
  title    = Tip,
  breakable,
}
\newtcolorbox{warnbox}{
  colback  = orange!8!white,
  colframe = orange!80!black,
  fonttitle= \bfseries,
  title    = Important,
  breakable,
}

\titleformat{\section}{\Large\bfseries}{\thesection}{0.7em}{}
\titleformat{\subsection}{\large\bfseries}{\thesubsection}{0.6em}{}
\titleformat{\subsubsection}{\normalsize\bfseries}{\thesubsubsection}{0.5em}{}

% ==============================================================================
\title{A Practical Tutorial for Magic VLSI Layout\\[4pt]
       \large Using the GlobalFoundries 180\,nm PDK (\texttt{gf180mcuD})}
\author{James E.\ Stine and Landon Burleson\\
        Oklahoma State University\\
        Electrical and Computer Engineering Department\\
        Stillwater, OK 74078 USA\\
        \texttt{james.stine@okstate.edu}}
\date{\today}
% ==============================================================================

\begin{document}
\maketitle
\tableofcontents
\newpage

% ==============================================================================
\section{Introduction}
% ==============================================================================

Magic is an open-source VLSI layout editor originally written by John
Ousterhout at UC~Berkeley and actively maintained by Tim Edwards at
\url{http://opencircuitdesign.com}.
It is the workhorse of the open-source analog/digital IC design flow,
providing interactive layout, built-in Design Rule Checking (DRC), and
netlist extraction---all in a single tool.
Magic receives frequent updates; always prefer version 8.3 or later.

This tutorial targets the \textbf{GlobalFoundries 180\,nm MCU PDK}
(\texttt{gf180mcuD}), an open-source process development kit available
through the \texttt{open\_pdks} installer and supported by Google's
open-silicon initiative.
The \texttt{gf180mcuD} process offers a 5-metal-layer CMOS flow 
suitable for both digital and mixed-signal designs.
All layer names, design rules, and command examples in this document
refer to \texttt{gf180mcuD} unless otherwise noted.

\begin{tipbox}
Always prepare a \textbf{stick diagram} before opening Magic.
Layout is a time sink; a hand sketch keeps you oriented and dramatically
reduces trial-and-error.
Think of it as your \emph{game plan}: spending just 5 minutes with a
stick diagram up front can save you hours of backtracking, erasing, and
re-routing inside Magic.
It forces you to think through transistor placement, wire routing, and
layer choices \emph{before} you are committed to geometry on screen---so
that when you do open Magic, you know exactly what to draw and in what
order.
Keep your stick diagrams in a lab notebook or a digital file you can
refer back to throughout the session.
\end{tipbox}

% ==============================================================================
\section{Prerequisites and Setup}
% ==============================================================================

\subsection{Software}

You need Magic (8.3+), \texttt{open\_pdks} with \texttt{gf180mcuD}
installed, and optionally KLayout for printing or viewing GDS output.
If you are using the IIC-OSIC-TOOLS Docker container
(\url{https://github.com/iic-jku/iic-osic-tools}), all of this is
pre-installed and pre-configured---skip ahead to
Section~\ref{sec:launching}.

For a native Linux build, follow the accompanying installation guide
and confirm that the following environment variables are set
(\texttt{setup.sh} handles this automatically):

\begin{lstlisting}
export OPEN_CIRCUITDESIGN_ROOT="$HOME/ocd"
export PDK_ROOT="$OPEN_CIRCUITDESIGN_ROOT/share/pdk"
export PDK=gf180mcuD
\end{lstlisting}

\subsection{Environment Check}

Before starting any design work, verify that Magic can see the PDK:

\begin{lstlisting}
ls $PDK_ROOT/gf180mcuD/libs.tech/magic/
# Should list: gf180mcuD.magicrc, gf180mcuD.tech, and related files
\end{lstlisting}

If the directory is empty or missing, re-run \texttt{open\_pdks} for
\texttt{gf180mcuD} or verify that \texttt{PDK\_ROOT} is set correctly.

\subsection{Create a Working Directory}

Organize each design in its own directory to avoid file collisions:

\begin{lstlisting}
mkdir -p ~/designs/inverter
cd ~/designs/inverter
\end{lstlisting}

% ==============================================================================
\section{Launching Magic}
\label{sec:launching}
% ==============================================================================

Always launch Magic with the \texttt{gf180mcuD} technology file so that
the correct DRC rules, layer colors, and extraction rules are loaded:

\begin{lstlisting}
magic -rcfile $PDK_ROOT/gf180mcuD/libs.tech/magic/gf180mcuD.magicrc
\end{lstlisting}

Magic opens two windows simultaneously:

\begin{itemize}[leftmargin=2em]
  \item \textbf{Graphics (layout) window} --- where you draw geometry.
  \item \textbf{Text (console) window} --- where commands are typed and
        all feedback, errors, and DRC messages appear.
\end{itemize}

\begin{warnbox}
Keep \emph{both} windows visible at all times.
DRC violations, extraction warnings, and all error messages appear
\emph{only} in the text window.
If something unexpected happens in the layout, look there first.
The graphical window will often appear completely normal even when serious problems exist --- a DRC violation does not flash or highlight itself automatically, and an extraction warning produces no visible artifact in the layout view. The \emph{only} indication that something has gone wrong is a message in the text window. Minimizing or hiding that window --- even briefly --- means you can miss critical feedback without realizing it.
A few practical habits follow from this:
\begin{itemize}
\item \textbf{After every major operation} (placing a cell, running DRC, extracting a net), glance at the text window before continuing.
\item \textbf{After a DRC run}, scroll through the full text output rather than stopping at the first clean result --- later violations may be more serious.
\item \textbf{When behavior seems wrong} --- a command appears to do nothing, a shape doesn't appear where expected, or a tool exits immediately --- the explanation is almost always already waiting in the text window.
\end{itemize}
Treat the text window not as a log to be consulted occasionally, but as a live diagnostic stream that runs in parallel with everything you do.
\end{warnbox}

% ==============================================================================
\section{The Magic Interface}
% ==============================================================================

\subsection{Mouse Buttons}

The mouse is your primary layout instrument.
Learn these three actions before anything else:

\begin{center}
\begin{tabular}{@{}llp{8cm}@{}}
\toprule
\textbf{Button} & \textbf{Shorthand} & \textbf{Action} \\
\midrule
Left click   & L & Sets the \emph{lower-left} corner of the selection box \\
Right click  & R & Sets the \emph{upper-right} corner of the selection box \\
Middle click & M & Paints the layer the cursor hovers over in the palette,
                   inside the current box \\
\bottomrule
\end{tabular}
\end{center}

\begin{tipbox}
The middle-click paint trick is one of Magic's fastest workflows: draw
your box with L and R, hover your cursor over the desired color in the
layer palette, then middle-click to fill immediately.
\end{tipbox}

\subsection{Entering Commands}

Commands can be entered in three equivalent ways:

\begin{enumerate}[leftmargin=2em]
  \item \textbf{Single-character macros} typed in the graphics window
        (e.g., \texttt{g} toggles the grid, \texttt{v} zooms to fit).
  \item \textbf{Long commands} prefixed with a colon typed in the
        \emph{text} window (e.g., \texttt{:paint poly},
        \texttt{:drc why}).
  \item \textbf{Mouse buttons} directly in the layout window for
        painting and selection.
\end{enumerate}

\subsection{The Box Tool}
All layout operations in Magic are relative to a rectangular
\emph{box}---the selection region shown as a white or yellow outline in
the layout window.
The default tool is the \textbf{BOX} tool, indicated by a crosshair
cursor.
To ensure you are in box mode:
\begin{lstlisting}
:tool box
\end{lstlisting}
The \texttt{Space} bar cycles through all tools (Box, Wiring, Netlist,
Pick).
For basic drawing, always remain in Box mode.
\clearpage

\begin{warnbox}
An accidental press of the \texttt{Space} bar will silently switch you
out of Box mode---and because the cursor change is subtle, it is sometimes
easy not to notice.
If something stops working as expected---painting a layer produces no
result, the box no longer resizes, or clicks seem to do the wrong
thing---your first check should be whether you have inadvertently cycled
to a different tool.
This is an extremely common source of confusion, especially when
typing commands while the cursor is over the layout window, since a
space in what you \emph{think} is the text window will instead be
captured by the layout window and advance the tool cycle.
To recover, either press \texttt{Space} repeatedly until the crosshair
cursor returns, or type the explicit command:
\begin{lstlisting}
:tool box
\end{lstlisting}
which will always return you to Box mode regardless of which tool is
currently active.
Make it a habit to run this command any time behavior feels
unexpected before attempting to diagnose anything else.
\end{warnbox}

% ==============================================================================
\section{Essential Command Reference}
% ==============================================================================

\subsection{File and View Commands}

\begin{center}
\begin{tabular}{@{}lp{9.5cm}@{}}
\toprule
\textbf{Command} & \textbf{Description} \\
\midrule
\texttt{:load \textit{name}}    & Load or create cell \textit{name}
                                  (creates a new empty cell if none exists) \\
\texttt{:save \textit{name}}    & Save current cell to
                                  \textit{name}\texttt{.mag} \\
\texttt{:quit}                  & Exit Magic \\
\texttt{:view} or \texttt{v}    & Zoom to fit all painted geometry \\
\texttt{:f}                     & Zoom to fit the entire cell boundary \\
\texttt{:zoom \textit{n}}       & Zoom in by factor $n$
                                  (\texttt{:zoom 0.5} zooms out) \\
\texttt{z} / \texttt{Z}         & Zoom window to box / zoom in $\times 2$ \\
\texttt{:grid} or \texttt{g}    & Toggle the visible grid \\
\texttt{:grid \textit{n}}       & Display a grid of
                                  $n\lambda \times n\lambda$ \\
\texttt{:grid off}              & Turn grid off \\
\texttt{:help [\textit{cmd}]}   & Print help for all or a specific command \\
\bottomrule
\end{tabular}
\end{center}

\subsection{Painting and Erasing}

\begin{lstlisting}
:paint <layer>      # paint <layer> inside the current box
:erase <layer>      # erase <layer> inside the current box
:erase *            # erase ALL layers inside the box
\end{lstlisting}

\subsection{Selection and Editing}

\begin{center}
\begin{tabular}{@{}lll@{}}
\toprule
\textbf{Command}     & \textbf{Macro} & \textbf{Description} \\
\midrule
\texttt{:select}     & \texttt{a} & Select geometry inside the box \\
\texttt{:copy}       & \texttt{c} & Copy the selection \\
\texttt{:move}       & ---        & Move the selection \\
\texttt{:delete}     & \texttt{d} & Delete the selection \\
\texttt{:undo}       & \texttt{u} & Undo last action \\
\texttt{:redo}       & \texttt{U} & Redo \\
\texttt{:sideways}   & ---        & Flip selection horizontally \\
\texttt{:upsidedown} & ---        & Flip selection vertically \\
\texttt{:clockwise}  & ---        & Rotate selection 90$^\circ$ clockwise \\
\bottomrule
\end{tabular}
\end{center}

Keys \texttt{Q W E R} (pressed in the graphics window) nudge the
selection by $1\lambda$ left, up, right, and down respectively.
The \texttt{.}\ (dot) macro repeats the last command---invaluable for
repetitive placement tasks.

% ==============================================================================
\section{gf180mcuD Layer Reference}
% ==============================================================================

The table below lists the most commonly used layers for basic cell
layout in \texttt{gf180mcuD}.
The \emph{Magic name} is exactly what you type after \texttt{:paint};
the description gives the physical meaning.

\begin{center}
\begin{tabular}{@{}llp{7cm}@{}}
\toprule
\textbf{Magic Name} & \textbf{Type} & \textbf{Description} \\
\midrule
\texttt{nwell}   & Implant  & N-type well (body of pMOS transistors) \\
\texttt{ndiff}   & Active   & N-type diffusion (nMOS source/drain) \\
\texttt{pdiff}   & Active   & P-type diffusion (pMOS source/drain) \\
\texttt{poly}    & Gate     & Polysilicon gate and short-range routing \\
\texttt{ndc}     & Contact  & N-type diffusion-to-metal contact \\
\texttt{pdc}     & Contact  & P-type diffusion-to-metal contact \\
\texttt{nsc}     & Contact  & N-type substrate contact (tie to GND) \\
\texttt{psc}     & Contact  & P-type substrate/well contact (tie to VDD) \\
\texttt{pc}      & Contact  & Poly-to-metal1 contact \\
\texttt{metal1}  & Metal    & Metal 1 --- primary routing layer \\
\texttt{metal2}  & Metal    & Metal 2 \\
\texttt{metal3}  & Metal    & Metal 3 \\
\texttt{metal4}  & Metal    & Metal 4 \\
\texttt{metal5}  & Metal    & Metal 5 (top layer) \\
\texttt{via1}    & Via      & metal1 $\rightarrow$ metal2 \\
\texttt{via2}    & Via      & metal2 $\rightarrow$ metal3 \\
\texttt{via3}    & Via      & metal3 $\rightarrow$ metal4 \\
\texttt{via4}    & Via      & metal4 $\rightarrow$ metal5 \\
\bottomrule
\end{tabular}
\end{center}

\begin{tipbox}
\textbf{Use metals for routing---not poly or local interconnect.}
Polysilicon and local interconnect (e.g., \texttt{li}) have
high resistance and capacitance.
These layers are more resistive
than \texttt{metal1} and should be reserved for short connections
where metal is physically impossible.
Route all signals on \texttt{metal1} and \texttt{metal2} whenever
possible.\\

Beyond resistance, there is a strong hierarchical reason to be
conservative with metals at the leaf-cell level: \textbf{stay at
metal1 and metal2 in leaf cells and leave the upper layers free.}
When you assemble leaf cells into blocks, and blocks into a full chip,
the higher metal layers are what the block-level and chip-level
routers depend on for power distribution, clock routing, and
inter-block signals.
If your leaf cells already consume metal3, metal4, or metal5, you
create conflicts that are painful---sometimes impossible---to resolve
without redesigning the cells from scratch.
Think of the metal stack as a shared resource across the entire
hierarchy: be a good citizen at the bottom, and the upper levels of
the design will be far easier to close.
\end{tipbox}

\subsection{Metal Stack Summary}

\texttt{gf180mcuD} provides five aluminum metal layers plus local
interconnect between poly and metal1:

\begin{center}
\begin{tabular}{@{}lll@{}}
\toprule
\textbf{Layer} & \textbf{Magic Name} & \textbf{Typical Use} \\
\midrule
Metal 1            & \texttt{metal1} & Leaf-cell routing (horizontal) \\
Metal 2            & \texttt{metal2} & Leaf-cell routing (vertical) \\
Metal 3            & \texttt{metal3} & Block-level routing \\
Metal 4            & \texttt{metal4} & Power distribution \\
Metal 5            & \texttt{metal5} & Top-level power/ground straps \\
\bottomrule
\end{tabular}
\end{center}

For leaf cells (individual gates and standard cells), try to stay at
metal1 and metal2 only.
Reserving upper metal layers for higher levels of hierarchy makes
block-level assembly and place-and-route much easier.
Alternate metal1 (horizontal) and metal2 (vertical) where possible to
minimize via counts and coupling.

% ==============================================================================
\section{Labeling and Ports}
% ==============================================================================

Labels attach logical names to layout geometry.
They are required for power and ground nodes and for any signal you
want to identify after netlist extraction.

\begin{lstlisting}
:label <name>
\end{lstlisting}

Position the box over the geometry to label, then run the command.
Follow these conventions:

\begin{itemize}[leftmargin=2em]
  \item \textbf{\texttt{vdd!}} and \textbf{\texttt{gnd!}} --- supply
        rails.
        The \texttt{!}\ suffix marks a node as \emph{global}, meaning
        it spans the entire design hierarchy.
        Type these exactly as shown.
  \item Signal names must match the schematic exactly; LVS compares
        them character-for-character.
  \item Adopt a consistent capitalization style within a project and
        keep it.
\end{itemize}

\begin{tipbox}
\textbf{Supply net naming and the global net convention.}
In Magic, appending \texttt{!} to a net name (e.g., \texttt{vdd!},
\texttt{gnd!}) is a convention that designates that net as
\emph{global}, meaning Magic will automatically connect every instance
of that label across the entire hierarchy without requiring an explicit
port or wire between cells.
This can be convenient for small designs, but it is \textbf{optional}
and entirely Magic-specific---it is not recognized by all tools, PDKs,
or LVS configurations.
Some flows define global supply nets through explicit port declarations,
named power domains, or technology-file directives instead, and in
those environments the \texttt{!} suffix may be ignored or behave
unexpectedly.\\
 
The more portable and robust habit is to label your supply wires with
consistent, simple names---\texttt{vdd} and \texttt{gnd}, or whatever
your schematic uses---and ensure those names match exactly between the
layout and the schematic netlist.
Whether or not you use the \texttt{!} convention, the rules that always
apply are:
\begin{itemize}
  \item \textbf{Be consistent} --- pick a naming convention and stick
        to it throughout the entire design.
  \item \textbf{Match the schematic} --- LVS compares net names
        character-for-character; \texttt{VDD} and \texttt{vdd} are
        different nets.
  \item \textbf{Label explicitly} --- a wire that is visually connected
        to a supply but carries no label will not be recognized as a
        supply net by the extractor.
\end{itemize}
\end{tipbox}

% ==============================================================================
\section{Step-by-Step: CMOS Inverter Layout in gf180mcuD}
% ==============================================================================

This section walks through drawing a minimum-size CMOS inverter from
scratch in \texttt{gf180mcuD}.
Have your stick diagram ready before you begin.

\subsection{Transistor Sizing}

In \texttt{gf180mcuD} the minimum drawn gate length is
$L_{\min} = 0.18\,\mu\text{m}$.
A reasonable starting inverter uses:

\begin{center}
\begin{tabular}{@{}llll@{}}
\toprule
\textbf{Transistor} & \textbf{Type} & \textbf{Width} & \textbf{Length} \\
\midrule
nMOS & N-channel & $0.42\,\mu\text{m}$ (minimum) & $0.18\,\mu\text{m}$ \\
pMOS & P-channel & $0.84\,\mu\text{m}$ ($2\times$ nMOS) & $0.18\,\mu\text{m}$ \\
\bottomrule
\end{tabular}
\end{center}

The pMOS is typically sized $2\times$ wider than the nMOS to equalize
rise and fall times, since hole mobility is roughly half that of
electrons.

\subsection{Step 1 --- Launch Magic and Create a New Cell}

\begin{lstlisting}
cd ~/designs/inverter
magic -rcfile $PDK_ROOT/gf180mcuD/libs.tech/magic/gf180mcuD.magicrc
\end{lstlisting}

In the text window, create a new cell named \texttt{inv}:

\begin{lstlisting}
:load inv
\end{lstlisting}

\subsection{Step 2 --- Draw the nMOS Diffusion}

Draw the nMOS active region in the p-substrate.
This defines the source and drain of the nMOS transistor.

\begin{lstlisting}
# nMOS source/drain region (in p-substrate)
:paint ndiff
\end{lstlisting}

\subsection{Step 3 --- Draw the pMOS Diffusion}

Draw the pMOS active region above (or beside) the nMOS region.
Leave enough space between them for the N-well enclosure and any
substrate/well taps you plan to add.

\begin{lstlisting}
# pMOS source/drain region
:paint pdiff
\end{lstlisting}

\subsection{Step 4 --- Draw the Polysilicon Gate}

The poly must cross \emph{both} active regions.
It forms the transistor gate wherever it overlaps diffusion.

\begin{lstlisting}
:paint poly
\end{lstlisting}

Draw one continuous vertical poly strip crossing both the
\texttt{ndiff} and \texttt{pdiff} regions.
The crossing automatically creates the gate of each transistor.

\begin{tipbox}
Keep dimensions at the DRC minimum where possible---smaller geometry
means less parasitic capacitance and faster switching.
That said, some layers are intentionally drawn wider than minimum:
supply rails, for example, are made wider to reduce resistance and
handle higher current without excessive voltage drop.
When in doubt, start using minimum geometries
and increase only when there is a
specific reason to do so.
\end{tipbox}

\subsection{Step 5 --- Add Diffusion Contacts}

Connect the diffusion source/drain regions up to metal1 using diffusion
contacts.

\begin{lstlisting}
# nMOS drain and source contacts
:paint ndc

# pMOS drain and source contacts
:paint pdc
\end{lstlisting}

Place contacts inside the source and drain regions of each transistor,
keeping them clear of the poly gate by the required spacing.

\subsection{Step 6 --- Add a Poly Contact}

To route the gate (poly) to a metal wire for the input signal:

\begin{lstlisting}
:paint pc
\end{lstlisting}

Place this at the top or bottom of the poly strip, wherever your input
wire will land.

\subsection{Step 7 --- Add Substrate and Well Taps}

Well and substrate taps are essential to prevent latch-up.
The p-substrate must be tied to \texttt{gnd!} and the N-well to
\texttt{vdd!}.

\begin{lstlisting}
# N-type substrate contact (p-substrate tie to GND)
:paint nsc

# P-type well contact (N-well tie to VDD)
:paint psc
\end{lstlisting}

Place \texttt{nsc} in the p-substrate and \texttt{psc} where the N-well
will be.
These taps also help define the final footprint of the pMOS region.

\subsection{Step 8 --- Route with Metal1}

Connect sources, drains, and gates using \texttt{metal1}:

\begin{lstlisting}
:paint metal1
\end{lstlisting}

Make the following connections:

\begin{itemize}[leftmargin=2em]
  \item nMOS source $\rightarrow$ \texttt{gnd!} (bottom power rail).
  \item pMOS source $\rightarrow$ \texttt{vdd!} (top power rail).
  \item nMOS drain + pMOS drain $\rightarrow$ output \texttt{Y}.
  \item nMOS gate + pMOS gate $\rightarrow$ input \texttt{A}.
\end{itemize}

\subsection{Step 9 --- Add Labels}

\begin{lstlisting}
# Position box over the VDD metal1 wire, then:
:label vdd!

# Position box over the GND metal1 wire, then:
:label gnd!

# Position box over the output metal1 wire, then:
:label Y

# Position box over the input wire (poly or metal), then:
:label A
\end{lstlisting}
\clearpage

\begin{tipbox}
When you need many similar labels---numbered pins, repeated bus
signals, or sequentially named wires---typing each one individually
in Magic can be tedious and error-prone.
Because Magic's layout files (\texttt{.mag}) are plain text, a useful
shortcut is to create one label, save the file, open it in any text
editor, and use find-and-replace to duplicate or rename the label
entries directly.
When you reopen the file in Magic, all of the edited labels appear
exactly as written---no retyping required.
This is especially handy for buses where only a trailing index number
changes (e.g., \texttt{data[0]} through \texttt{data[7]}).
\end{tipbox}

\subsection{Step 10 --- Draw the N-Well}

Now that the pMOS diffusion, well tap, and metal routing are all in
place, you can see exactly how much area the pMOS region occupies.
Draw the N-well last so you can size it precisely to enclose the
\texttt{pdiff} and \texttt{psc} regions with the correct enclosure
margin---no guesswork required.

\begin{lstlisting}
:paint nwell
\end{lstlisting}

Extend the N-well until it fully covers the \texttt{pdiff} and
\texttt{psc} geometry with sufficient overlap on all sides.
DRC will immediately highlight any enclosure violation if the N-well is
too tight.

\subsection{Step 11 --- Run DRC}

\begin{lstlisting}
:drc check
:drc why
\end{lstlisting}

Magic runs DRC continuously in the background; \texttt{:drc why}
explains any violation under or near the current box.
Fix all violations before proceeding to extraction.

\begin{warnbox}
Do not proceed to extraction or LVS until DRC reports zero errors.
A clean DRC is a hard prerequisite for a meaningful LVS comparison.
\end{warnbox}

\subsection{Step 12 --- Save}

\begin{lstlisting}
:save inv
\end{lstlisting}

This writes \texttt{inv.mag} to your current working directory.
Save frequently; Magic does not auto-save.

% ==============================================================================
\section{Design Rule Checking (DRC)}
% ==============================================================================

Magic performs DRC continuously as you draw, highlighting violations in
real time.
The following commands are most useful during layout:

\begin{center}
\begin{tabular}{@{}lp{9.5cm}@{}}
\toprule
\textbf{Command}    & \textbf{Description} \\
\midrule
\texttt{:drc check} & Force a complete DRC pass \\
\texttt{:drc why}   & Explain the DRC error inside the current box \\
\texttt{:drc find}  & Move the box to the next DRC error \\
\texttt{:drc count} & Report the total number of DRC errors \\
\bottomrule
\end{tabular}
\end{center}

Common DRC errors in \texttt{gf180mcuD} and how to fix them:

\begin{itemize}[leftmargin=2em]
  \item \textbf{Poly too close to diffusion edge} --- increase the
        spacing between the poly gate end and the edge of the diffusion
        region.
  \item \textbf{Contact enclosure violation} --- the surrounding
        diffusion or metal region is too narrow; enlarge the enclosing
        layer around the contact.
  \item \textbf{N-well enclosure of pdiff insufficient} --- the N-well
        must extend beyond the p-type diffusion by the required
        enclosure distance; widen the N-well.
  \item \textbf{Metal spacing violation} --- two metal1 wires are too
        close together; increase the gap between them.
  \item \textbf{Via overhang violation} --- vias in \texttt{gf180mcuD}
        require a specific enclosure overhang on both the lower and
        upper metal layers; ensure the metal extends far enough beyond
        the via on all four sides.
\end{itemize}

% ==============================================================================
\section{Extraction and Netlist Generation}
% ==============================================================================

Once DRC is clean, extract a netlist from the layout for LVS and
post-layout simulation.

\begin{lstlisting}
:extract all
:ext2spice hierarchy on
:ext2spice format ngspice
:ext2spice
\end{lstlisting}

This produces \texttt{inv.ext} (the intermediate extraction database)
and \texttt{inv.spice} (a SPICE netlist with parasitic capacitances).
You can also use the course convenience script if it is available in
your environment:

\begin{lstlisting}
ext4mag inv.mag
\end{lstlisting}

This script generates both \texttt{.gds} and \texttt{.sim} files and
optionally a SPICE deck (answer \texttt{N} when prompted about SPICE if
you do not need it yet).

% ==============================================================================
\section{Layout vs.\ Schematic (LVS)}
% ==============================================================================

LVS verifies that your layout correctly implements the intended
schematic.
It is one of the most critical steps in the design flow.
Many professional organizations list skipping LVS as a termination
offense.
Run LVS after every cell passes DRC, and use it as a progressive
debugging tool as you assemble larger designs cell by cell.

Run LVS with Netgen:

\begin{lstlisting}
netgen -batch lvs "inv.spice inv" \
    "$PDK_ROOT/gf180mcuD/libs.tech/netgen/gf180mcuD_setup.tcl"
\end{lstlisting}

A passing result prints \textbf{Circuits match} with zero
discrepancies.
If it fails, Netgen lists the mismatched devices or nets---trace these
back to errors in your layout or labels.

\begin{tipbox}
Most companies draw the schematic first, then create the layout to
match it.
A SPICE netlist from xschem or a textual HDL description both serve as
valid references for Netgen.
\end{tipbox}

% ==============================================================================
\section{Plotting and Printing Your Layout}
% ==============================================================================

Always use KLayout to produce layout documentation---never screenshots.
KLayout renders GDS layers accurately and exports to postscript or PDF.

\begin{lstlisting}
# Inside Magic: write GDS
:gds write inv

# In a terminal: open with KLayout
klayout inv.gds
\end{lstlisting}

To convert KLayout's postscript output to PDF:

\begin{lstlisting}
ps2pdf ./inv.ps inv.pdf
\end{lstlisting}

KLayout should be installed as part of your EDA tool suite.
If layer colors do not appear correctly, import the
\texttt{gf180mcuD} technology file via
\textbf{Tools $\rightarrow$ Manage Technologies $\rightarrow$
Import Technology}.

% ==============================================================================
\section{Switch-Level Simulation with IRSIM}
% ==============================================================================

After a successful LVS, validate digital logic function with IRSIM, a
switch-level simulator that models transistors as ideal switches.
IRSIM is much faster than full SPICE simulation and is well-suited for
verifying logic correctness from extracted layout.

\subsection{Extraction for IRSIM}

\begin{lstlisting}
# In Magic
:extract all
:ext2sim
\end{lstlisting}

This produces \texttt{inv.sim}, which IRSIM reads directly.

\subsection{Running IRSIM}

\begin{lstlisting}
irsim $PDK_ROOT/gf180mcuD/libs.tech/irsim/gf180mcuD.prm inv.sim
\end{lstlisting}

\subsection{IRSIM Command Files}

Automate stimulus with a \texttt{.cmd} file to build a reusable test
vector set (often called a \emph{golden file}).
A simple inverter command file:

\begin{lstlisting}
# inv.cmd
stepsize 100
| Initial state: A = 0, expect Y = 1
set A 0
c
print Y
| Toggle: A = 1, expect Y = 0
set A 1
c
print Y
| Toggle back: A = 0, expect Y = 1
set A 0
c
print Y
\end{lstlisting}

Load and run the command file inside IRSIM:

\begin{lstlisting}
source inv.cmd
\end{lstlisting}

A correct inverter outputs \texttt{Y} high when \texttt{A} is low, and
\texttt{Y} low when \texttt{A} is high.

% ==============================================================================
\section{The Complete Design Flow}
% ==============================================================================

The full open-source IC design flow for \texttt{gf180mcuD} is:

\begin{enumerate}[leftmargin=2em]
  \item \textbf{Draw schematic} in xschem or write a SPICE text netlist.
  \item \textbf{Simulate schematic} with ngspice (SPICE) or IRSIM for
        initial functional verification.
  \item \textbf{Draw stick diagram} to plan transistor placement and
        wire routing.
  \item \textbf{Create layout} in Magic using the \texttt{gf180mcuD}
        technology.
  \item \textbf{Run DRC} --- fix all violations before moving forward.
  \item \textbf{Run LVS} with Netgen against the schematic netlist.
  \item \textbf{Extract and simulate layout} --- run ngspice or IRSIM
        on the extracted netlist to capture parasitics and verify timing
        and power.
  \item \textbf{Plot} the layout using KLayout for documentation and
        submission.
\end{enumerate}
\clearpage

\begin{tcolorbox}[
  colback  = red!8!white,
  colframe = red!70!black,
  fonttitle= \bfseries,
  title    = Important,
  breakable,
]
\textbf{Never skip LVS. This is not optional.}
LVS is your guarantee that what you drew is actually what you designed.
Skipping it is one of the most costly mistakes you can make in IC
design---a layout that passes DRC but fails LVS will produce a chip that
does not work, and at fabrication costs of thousands to hundreds of
thousands of dollars per run, there is no opportunity to just ``try
again.''
The consequences in industry are severe: I have heard from colleagues
that several companies explicitly
list bypassing LVS as grounds for immediate dismissal or a reprimand.
Run LVS after every cell passes DRC---not just at the top level, not
just at the end of the project, but cell by cell as you build up the
design.
Catching a net mismatch in a single gate takes minutes to fix; finding
the same error buried in a full-chip LVS failure can take days.\\
 
Be warned: \textbf{getting LVS to pass is genuinely extremely hard.}
A failing LVS report is not a simple error message---it is a puzzle.
The tool gives you clues in the form of mismatched nets, unmatched
devices, and node-count discrepancies, and it is up to you to decode
what they mean and trace them back to the root cause in your layout or
labels.
Approach it like a debugging session: read the output carefully, form a
hypothesis, make one change, and re-run. \\
 
Commercial tools such as Mentor/Siemens Calibre~DRV are the
industry gold standard for this, offering rich GUI integrations that
highlight mismatches directly in the layout and schematic side by side,
making the puzzle significantly easier to solve.
Open-source engines such as Netgen are capable and free, but they are
typically slower and produce less guided output than their commercial
counterparts.
That is not a reason to avoid them---it is a reason to stay
vigilant and optimistic.
Read the output carefully, work through it methodically, and remember
that a passing LVS is always achievable if your layout is correct.
\end{tcolorbox}

% ==============================================================================
\section{Tips and Common Pitfalls}
% ==============================================================================

\begin{itemize}[leftmargin=2em]
  \item \textbf{Always use a stick diagram.}
        The single most effective way to avoid wasting hours in Magic is
        to know exactly what you intend to draw before you open the
        tool.

  \item \textbf{Route on metal, not poly.}
        Poly has significantly higher resistance and capacitance than
        metal.
        Use it only for the transistor gate and very short local
        connections.

  \item \textbf{Stay at metal1/metal2 in leaf cells.}
        Reserve upper metal layers for block-level routing and power.
        Wires are the dominant contributor to cell area; keeping them
        at the lowest layers preserves routing resources for higher
        levels of hierarchy.

  \item \textbf{Alternate metal directions.}
        Use metal1 for horizontal routing and metal2 for vertical
        routing (HVH style) to minimize via counts and reduce coupling.

  \item \textbf{Use pitch as a design-wide parameter.}
        Establish a consistent wire pitch for all wires in a cell and
        stick to it.
        Uniform pitch makes cell abutment and later place-and-route
        much smoother.

  \item \textbf{Label everything.}
        Supply rails (\texttt{vdd!}, \texttt{gnd!}) and all I/O
        signals must be labeled for extraction and LVS to work
        correctly.

  \item \textbf{Save frequently.}
        Use \texttt{:save} after any significant change.
        Magic does not auto-save.

  \item \textbf{Use scripts.}
        Extraction, LVS, and GDS export are all scriptable.
        A well-written script saves hours across a semester of design
        work.
        You can find the location of any installed script with
        \texttt{which \textit{scriptname}}.

  \item \textbf{Post problems publicly.}
        Include a screenshot and a step-by-step description of what
        you tried.
        Other students learn from your problems, and public posts build
        a searchable archive of solutions for the entire class.
\end{itemize}

% ==============================================================================
\section{Quick-Reference Cheat Sheet}
% ==============================================================================

\subsection*{Keyboard Macros}

\begin{center}
\begin{tabular}{@{}lll@{}}
\toprule
\textbf{Key} & \textbf{Long form} & \textbf{Action} \\
\midrule
\texttt{v}       & \texttt{:view}     & Zoom to fit all geometry \\
\texttt{g}       & \texttt{:grid}     & Toggle grid \\
\texttt{a}       & \texttt{:select}   & Select inside box \\
\texttt{c}       & \texttt{:copy}     & Copy selection \\
\texttt{d}       & \texttt{:delete}   & Delete selection \\
\texttt{u}       & \texttt{:undo}     & Undo \\
\texttt{U}       & \texttt{:redo}     & Redo \\
\texttt{z}       & ---                & Zoom window to box \\
\texttt{Z}       & \texttt{:zoom 2}   & Zoom in $\times 2$ \\
\texttt{.}       & ---                & Repeat last command \\
\texttt{Q/W/E/R} & ---                & Nudge selection $1\lambda$
                                        left/up/right/down \\
Space            & \texttt{:tool box} & Cycle tools (return to Box) \\
\bottomrule
\end{tabular}
\end{center}

\subsection*{Key Commands}

\noindent To launch Magic, use (where \texttt{\$PDK\_ROOT} is set by \texttt{setup.sh}):
\begin{lstlisting}
magic -rcfile $PDK_ROOT/gf180mcuD/libs.tech/magic/gf180mcuD.magicrc
\end{lstlisting}

\noindent Alternatively, you can omit the \texttt{-rcfile} flag by placing a
\texttt{.magicrc} file in your home or working directory that loads the
technology automatically. Add the following line to \texttt{\textasciitilde/.magicrc}:
\begin{lstlisting}
tech load $env(PDK_ROOT)/gf180mcuD/libs.tech/magic/gf180mcuD.tech
\end{lstlisting}
\noindent With this in place, simply typing \texttt{magic} from any directory
will automatically load the \texttt{gf180mcuD} technology without requiring
the \texttt{-rcfile} argument each time.

\begin{center}
\begin{tabular}{@{}lp{7.5cm}@{}}
\toprule
\textbf{Command} & \textbf{Purpose} \\
\midrule
\texttt{:drc why}        & Explain DRC error inside current box \\
\texttt{:drc find}       & Jump box to next DRC error \\
\texttt{:extract all}    & Extract netlist from layout \\
\texttt{:ext2spice}      & Write SPICE from extraction database \\
\texttt{:ext2sim}        & Write .sim file for IRSIM \\
\texttt{:gds write inv}  & Export GDS file \\
\texttt{klayout inv.gds} & View or print GDS with KLayout \\
\bottomrule
\end{tabular}
\end{center}

\subsection*{Common Paint Commands for gf180mcuD}

\begin{center}
\begin{tabular}{@{}ll@{}}
\toprule
\textbf{Command} & \textbf{What it draws} \\
\midrule
\texttt{:paint nwell}  & N-well (pMOS body) \\
\texttt{:paint ndiff}  & N-type diffusion \\
\texttt{:paint pdiff}  & P-type diffusion \\
\texttt{:paint poly}   & Polysilicon gate \\
\texttt{:paint ndc}    & N diffusion contact \\
\texttt{:paint pdc}    & P diffusion contact \\
\texttt{:paint nsc}    & N substrate contact \\
\texttt{:paint psc}    & P well contact \\
\texttt{:paint pc}     & Poly contact \\
\texttt{:paint metal1} & Metal 1 \\
\texttt{:paint metal2} & Metal 2 \\
\texttt{:paint via1}   & Via between metal1 and metal2 \\
\bottomrule
\end{tabular}
\end{center}

% ==============================================================================
\section*{Additional Resources}
\addcontentsline{toc}{section}{Additional Resources}
% ==============================================================================

\begin{itemize}[leftmargin=2em]
  \item Magic documentation and tutorials:
        \url{http://opencircuitdesign.com/magic}
  \item gf180mcuD PDK documentation:
        \url{https://gf180mcu-pdk.readthedocs.io}
  \item open\_pdks repository (PDK installer):
        \url{https://github.com/RTimothyEdwards/open_pdks}
  \item IIC-OSIC-TOOLS Docker container:
        \url{https://github.com/iic-jku/iic-osic-tools}
  \item Course examples, scripts, and layout guides:
        \url{https://stineje.github.io}
  \item Video walkthroughs:
        \url{https://www.youtube.com/user/jlstine/}
\end{itemize}

\end{document}
